\documentclass{article}
\usepackage{amsmath,amssymb,amsthm}
\usepackage{algorithm}
\usepackage{algpseudocode}
\usepackage{geometry}
\geometry{margin=1in}

\newtheorem{theorem}{Theorem}
\newtheorem{lemma}{Lemma}
\newtheorem{definition}{Definition}
\newtheorem{assumption}{Assumption}
\newtheorem{corollary}{Corollary}

\begin{document}

\section*{Randomized Coordinate Gradient Descent (RCGD) under PL Condition}

We analyze a \emph{Randomized Coordinate Gradient Descent} method designed for
functions $f:\mathbb{R}^d \to \mathbb{R}$ that are smooth coordinate-wise and
satisfy the Polyak--\L ojasiewicz (PL) inequality. At each iteration, a single
coordinate $i_t \in \{1,\dots,d\}$ is sampled uniformly at random, the partial
derivative $\nabla_{i_t} f(w_t)$ is computed, and the parameter vector is updated
only along that coordinate.

\section*{Mathematical Formulation}

Let $f:\mathbb{R}^d \to \mathbb{R}$ and let $e_i \in \mathbb{R}^d$ denote the
$i$-th standard basis vector. Define the partial derivative
\[
\nabla_i f(w) := \frac{\partial f(w)}{\partial w_i}, \qquad
\nabla f(w) = \big(\nabla_1 f(w), \dots, \nabla_d f(w)\big)^\top .
\]

\noindent\textbf{Update rule.} At iteration $t$:
\begin{enumerate}
\item Sample $i_t \sim \mathrm{Uniform}\{1,\dots,d\}$ independently.
\item Compute the coordinate gradient $\nabla_{i_t} f(w_t)$.
\item Update only coordinate $i_t$:
\[
w_{t+1} = w_t - \eta\, \nabla_{i_t} f(w_t)\, e_{i_t}.
\]
\end{enumerate}

\noindent\textbf{Hyperparameters.}
\begin{itemize}
\item Step size $\eta > 0$, typically chosen as $\eta = 1/L$ where $L$ is the
coordinate-wise smoothness constant.
\item Number of iterations $T$ (or a gradient tolerance $\epsilon$).
\end{itemize}

\noindent\textbf{Componentwise form.} For $j = 1,\dots,d$,
\[
(w_{t+1})_j =
\begin{cases}
(w_t)_j - \eta\, \nabla_j f(w_t), & j = i_t,\\[4pt]
(w_t)_j, & j \neq i_t.
\end{cases}
\]

\section*{Pseudocode}

\begin{algorithm}[H]
\caption{Randomized Coordinate Gradient Descent (RCGD)}
\begin{algorithmic}[1]
\Require Initial point $w_0 \in \mathbb{R}^d$, step size $\eta > 0$, iteration budget $T$, tolerance $\epsilon$
\State $t \gets 0$
\While{$t < T$ \textbf{and} $\|\nabla f(w_t)\| > \epsilon$}
    \State Sample $i_t \sim \mathrm{Uniform}\{1,\dots,d\}$
    \State Compute partial derivative $g \gets \nabla_{i_t} f(w_t)$
    \State $w_{t+1} \gets w_t - \eta\, g\, e_{i_t}$ \Comment{update only coordinate $i_t$}
    \State $t \gets t + 1$
\EndWhile
\State \Return $w_t$
\end{algorithmic}
\end{algorithm}

\section*{Dry Run Example}

Consider $f(w) = (w-3)^2$, a one-dimensional problem ($d=1$), with
$w_0 = 0$, $\eta = 0.1$. Since $d=1$, the only coordinate is selected at every
step, and $\nabla f(w) = 2(w-3)$.

\begin{itemize}
\item \textbf{Iteration 1:}
\[
\nabla f(w_0) = 2(0-3) = -6, \quad
w_1 = 0 - 0.1\cdot(-6) = 0.6, \quad
f(w_1) = (0.6-3)^2 = 5.76.
\]
\item \textbf{Iteration 2:}
\[
\nabla f(w_1) = 2(0.6-3) = -4.8, \quad
w_2 = 0.6 - 0.1\cdot(-4.8) = 1.08, \quad
f(w_2) = (1.08-3)^2 = 3.6864.
\]
\item \textbf{Iteration 3:}
\[
\nabla f(w_2) = 2(1.08-3) = -3.84, \quad
w_3 = 1.08 - 0.1\cdot(-3.84) = 1.464, \quad
f(w_3) = (1.464-3)^2 \approx 2.3593.
\]
\end{itemize}

The objective decreases monotonically: $9 \to 5.76 \to 3.6864 \to 2.3593$,
converging toward the optimum $w^* = 3$, $f(w^*)=0$.

\section*{Assumptions}

\begin{assumption}[Coordinate-wise smoothness]
\label{ass:smooth}
There exists $L > 0$ such that for every $i \in \{1,\dots,d\}$, every $w \in \mathbb{R}^d$,
and every $\alpha \in \mathbb{R}$,
\[
\left| \nabla_i f(w + \alpha e_i) - \nabla_i f(w) \right| \le L\, |\alpha|.
\]
Equivalently, the descent lemma along coordinate $i$ holds:
\[
f(w + \alpha e_i) \le f(w) + \alpha\, \nabla_i f(w) + \frac{L}{2}\alpha^2 .
\]
\end{assumption}

\begin{assumption}[Polyak--\L ojasiewicz (PL) condition]
\label{ass:pl}
There exists $\mu > 0$ such that for all $w \in \mathbb{R}^d$,
\[
\frac{1}{2}\|\nabla f(w)\|^2 \ge \mu\big(f(w) - f^*\big),
\]
where $f^* = \inf_w f(w) > -\infty$.
\end{assumption}

\begin{assumption}[Lower boundedness]
\label{ass:lb}
$f$ is bounded below: $f^* = \inf_{w} f(w) > -\infty$, attained or approached.
\end{assumption}

\begin{assumption}[Step size constraint]
\label{ass:step}
The step size satisfies $0 < \eta \le \dfrac{1}{L}$.
\end{assumption}

\section*{Main Convergence Theorem}

\begin{theorem}[Linear convergence of RCGD under PL]
\label{thm:main}
Under Assumptions~\ref{ass:smooth}--\ref{ass:step}, with step size
$\eta = 1/L$, the iterates of RCGD satisfy, for all $t \ge 0$,
\[
\mathbb{E}\big[f(w_t) - f^*\big]
\le \left(1 - \frac{\mu}{dL}\right)^{t} \big(f(w_0) - f^*\big).
\]
Equivalently, to reach $\mathbb{E}[f(w_T)-f^*] \le \epsilon$, it suffices that
\[
T \ge \frac{dL}{\mu}\,\log\!\frac{f(w_0)-f^*}{\epsilon}.
\]
\end{theorem}

\section*{Theoretical Properties}

\begin{itemize}
\item \textbf{Linear (geometric) rate.} The contraction factor
$\rho = 1 - \mu/(dL) \in (0,1)$ guarantees a linear rate, despite no convexity
assumption (PL is weaker than strong convexity).
\item \textbf{Dimension dependence.} The factor $d$ reflects that each iteration
touches only one of $d$ coordinates; full-gradient GD has rate $1-\mu/L$ but
costs $d$ times more per iteration. The \emph{per-coordinate-evaluation}
complexity is comparable.
\item \textbf{Hyperparameter sensitivity.} The rate is optimized at
$\eta = 1/L$. Underestimating $L$ (too-large $\eta$) violates
Assumption~\ref{ass:step} and may break the descent guarantee; overestimating
$L$ slows convergence proportionally.
\item \textbf{Stability.} Because each update applies the descent lemma along a
single coordinate, the objective decreases in expectation at every step
(supermartingale-like behaviour of $f(w_t)-f^*$).
\end{itemize}

\section*{Discussion}

RCGD is attractive when partial derivatives are much cheaper than full
gradients (e.g., sparse/separable structure). Under the PL condition it attains
linear convergence to the global optimal value $f^*$ without requiring
uniqueness of minimizers or convexity, generalizing classical coordinate-descent
analyses. The randomized coordinate selection makes the expected progress
isotropic, yielding the clean $1-\mu/(dL)$ factor.

\section*{Preliminaries (Definitions and Lemmas)}

We work with the natural filtration
$\mathcal{F}_t = \sigma(i_0,\dots,i_{t-1})$, so that $w_t$ is
$\mathcal{F}_t$-measurable and $i_t$ is independent of $\mathcal{F}_t$.

\begin{lemma}[Coordinate descent lemma]
\label{lem:descent}
Under Assumption~\ref{ass:smooth}, for any $w$, any coordinate $i$, and update
$w^+ = w - \eta\, \nabla_i f(w)\, e_i$,
\[
f(w^+) \le f(w) - \eta\left(1 - \frac{L\eta}{2}\right)\big(\nabla_i f(w)\big)^2 .
\]
\end{lemma}

\begin{proof}
Apply Assumption~\ref{ass:smooth} with $\alpha = -\eta\,\nabla_i f(w)$:
\[
f(w + \alpha e_i) \le f(w) + \alpha\,\nabla_i f(w) + \frac{L}{2}\alpha^2 .
\]
Substituting $\alpha = -\eta\,\nabla_i f(w)$:
\[
f(w^+) \le f(w) - \eta\big(\nabla_i f(w)\big)^2 + \frac{L}{2}\eta^2 \big(\nabla_i f(w)\big)^2
= f(w) - \eta\left(1 - \frac{L\eta}{2}\right)\big(\nabla_i f(w)\big)^2 .
\]
\end{proof}

\section*{Main Proof (Step-by-Step Derivation)}

\begin{proof}[Proof of Theorem~\ref{thm:main}]

\textbf{[Step 1] Apply the coordinate descent lemma.}
Conditioned on $\mathcal{F}_t$ and on the random index $i_t$, the iterate is
$w_{t+1} = w_t - \eta\,\nabla_{i_t} f(w_t)\,e_{i_t}$. By
Lemma~\ref{lem:descent},
\[
f(w_{t+1}) \le f(w_t) - \eta\left(1 - \frac{L\eta}{2}\right)\big(\nabla_{i_t} f(w_t)\big)^2 .
\]

\textbf{[Step 2] Take conditional expectation over $i_t$.}
Since $i_t$ is uniform on $\{1,\dots,d\}$ and independent of $\mathcal{F}_t$,
\[
\mathbb{E}\big[(\nabla_{i_t} f(w_t))^2 \mid \mathcal{F}_t\big]
= \sum_{i=1}^{d} \frac{1}{d}\big(\nabla_i f(w_t)\big)^2
= \frac{1}{d}\sum_{i=1}^{d}\big(\nabla_i f(w_t)\big)^2
= \frac{1}{d}\,\|\nabla f(w_t)\|^2 .
\]
Taking conditional expectation of the Step 1 inequality (the term $f(w_t)$ is
$\mathcal{F}_t$-measurable):
\[
\mathbb{E}\big[f(w_{t+1}) \mid \mathcal{F}_t\big]
\le f(w_t) - \frac{\eta}{d}\left(1 - \frac{L\eta}{2}\right)\|\nabla f(w_t)\|^2 .
\]

\textbf{[Step 3] Substitute the step size $\eta = 1/L$.}
Then $1 - \tfrac{L\eta}{2} = 1 - \tfrac{1}{2} = \tfrac{1}{2}$ and
$\tfrac{\eta}{d} = \tfrac{1}{dL}$, giving
\[
\mathbb{E}\big[f(w_{t+1}) \mid \mathcal{F}_t\big]
\le f(w_t) - \frac{1}{2dL}\|\nabla f(w_t)\|^2 .
\]

\textbf{[Step 4] Apply the PL inequality.}
By Assumption~\ref{ass:pl}, $\tfrac{1}{2}\|\nabla f(w_t)\|^2 \ge \mu(f(w_t)-f^*)$,
i.e. $\|\nabla f(w_t)\|^2 \ge 2\mu\,(f(w_t)-f^*)$. Substituting (the coefficient
$\tfrac{1}{2dL}>0$, so the inequality direction is preserved):
\[
\mathbb{E}\big[f(w_{t+1}) \mid \mathcal{F}_t\big]
\le f(w_t) - \frac{1}{2dL}\cdot 2\mu\,(f(w_t)-f^*)
= f(w_t) - \frac{\mu}{dL}\,(f(w_t)-f^*).
\]

\textbf{[Step 5] Subtract $f^*$ and rearrange.}
Subtract $f^*$ from both sides:
\[
\mathbb{E}\big[f(w_{t+1}) - f^* \mid \mathcal{F}_t\big]
\le (f(w_t)-f^*) - \frac{\mu}{dL}(f(w_t)-f^*)
= \left(1 - \frac{\mu}{dL}\right)(f(w_t)-f^*).
\]

\textbf{[Step 6] Take total expectation (tower property).}
Applying $\mathbb{E}[\cdot] = \mathbb{E}[\mathbb{E}[\cdot \mid \mathcal{F}_t]]$,
\[
\mathbb{E}\big[f(w_{t+1}) - f^*\big]
\le \left(1 - \frac{\mu}{dL}\right)\mathbb{E}\big[f(w_t)-f^*\big].
\]

\textbf{[Step 7] Unroll the recursion.}
Let $\delta_t := \mathbb{E}[f(w_t)-f^*]$ and $\rho := 1 - \tfrac{\mu}{dL}$.
Step 6 reads $\delta_{t+1} \le \rho\,\delta_t$. Iterating from $0$ to $t$,
\[
\delta_t \le \rho^t \delta_0,
\qquad\text{i.e.}\qquad
\mathbb{E}\big[f(w_t)-f^*\big] \le \left(1 - \frac{\mu}{dL}\right)^t (f(w_0)-f^*).
\]

\textbf{[Step 8] Verify $\rho \in (0,1)$.}
By PL combined with $L$-smoothness one has $\mu \le L \le dL$, hence
$0 < \mu/(dL) \le 1$ and $\rho = 1-\mu/(dL) \in [0,1)$, so the bound contracts.
\end{proof}

\section*{Rate Derivation (With Constants)}

To guarantee $\mathbb{E}[f(w_T)-f^*] \le \epsilon$, using
$\delta_T \le \rho^T \delta_0$ with $\rho = 1-\mu/(dL)$, we require
\[
\left(1 - \frac{\mu}{dL}\right)^{T}(f(w_0)-f^*) \le \epsilon .
\]
Taking logarithms and using $\log(1-x) \le -x$ for $x \in [0,1)$,
\[
T\,\log\!\left(1-\frac{\mu}{dL}\right) \le \log\frac{\epsilon}{f(w_0)-f^*}
\;\Longleftarrow\;
-T\,\frac{\mu}{dL} \le \log\frac{\epsilon}{f(w_0)-f^*},
\]
which yields the sufficient condition
\[
T \ge \frac{dL}{\mu}\,\log\!\frac{f(w_0)-f^*}{\epsilon}
= O\!\left(\frac{dL}{\mu}\log\frac{1}{\epsilon}\right).
\]

\section*{Special Cases}

\paragraph{General step size $\eta \le 1/L$.}
From Step 2, keeping $\eta$ general and applying PL,
\[
\mathbb{E}[f(w_{t+1})-f^*]
\le \left(1 - \frac{2\mu\eta}{d}\Big(1-\frac{L\eta}{2}\Big)\right)
\mathbb{E}[f(w_t)-f^*].
\]
The contraction factor $g(\eta)=\tfrac{2\mu\eta}{d}\big(1-\tfrac{L\eta}{2}\big)$
is maximized by $g'(\eta)=0 \Rightarrow 1-L\eta = 0 \Rightarrow \eta = 1/L$,
recovering the optimal factor $\mu/(dL)$ of Theorem~\ref{thm:main}.

\paragraph{Strongly convex $f$.}
If $f$ is $\mu$-strongly convex, the PL inequality holds with the same $\mu$,
so Theorem~\ref{thm:main} applies verbatim, recovering the classical linear rate
$\big(1-\mu/(dL)\big)^t$ of randomized coordinate descent.

\paragraph{Dimension $d=1$ (full gradient).}
The method reduces to gradient descent and the rate becomes
$\big(1-\mu/L\big)^t$, matching the standard PL convergence rate for GD; the
dry-run example above is an instance of this case with monotone decrease.

\paragraph{Gradient-norm guarantee.}
Combining Step 3 with the linear rate, summing $\|\nabla f(w_t)\|^2$ over
$t=0,\dots,T-1$ gives
\[
\frac{1}{T}\sum_{t=0}^{T-1}\mathbb{E}\|\nabla f(w_t)\|^2
\le \frac{2dL\,(f(w_0)-f^*)}{T} = O\!\left(\frac{dL}{T}\right),
\]
showing an $O(1/T)$ stationarity rate even without invoking PL.

\end{document}